% !TITLE 鹊桥仙·纤云弄巧
% !AUTHOR 秦观
% !DYNASTY 两宋
% !GENRE 词
% !ORDER 10

\begin{poem}
纤云弄巧\textsuperscript{1}，飞星传恨\textsuperscript{2}，银汉\textsuperscript{3}迢迢暗度。
金风玉露\textsuperscript{4}一相逢，便胜却人间无数。

柔情似水，佳期如梦，忍顾\textsuperscript{5}鹊桥\textsuperscript{6}归路。
两情若是久长时，又岂在朝朝暮暮\textsuperscript{7}。
\end{poem}

\begin{annotations}
\notetext{1}{纤云弄巧：纤细的云彩变化出巧妙的形状。七夕的主题是乞巧（女孩子向织女祈求灵巧的手艺），“弄巧”暗扣七夕。}
\notetext{2}{飞星传恨：流星传递着离恨。飞星指划过天空的流星，也暗指牛郎织女星。}
\notetext{3}{银汉：银河、天河。牛郎织女被银河阻隔，七夕夜才能暗暗渡过天河相会。}
\notetext{4}{金风玉露：秋风和白露。金风，秋天五行属金，故称秋风为金风；玉露，晶莹如玉的露水。在秋风白露的季节里相逢一次，胜过人间无数次的相会。}
\notetext{5}{忍顾：不忍回头看。怎么忍心回头看那座鹊桥上的归路呢——相聚太短，分离在即。}
\notetext{6}{鹊桥：喜鹊搭成的桥。传说七夕夜喜鹊们在银河上搭桥，让牛郎织女过河相会。}
\notetext{7}{朝朝暮暮：每天每夜。只要两情长久，何必要天天厮守呢——这句成了中国爱情观中最洒脱也是最大度的表达。}
\end{annotations}
